\scanpage{7}
% L1–5
et de $\bigwedge_{i\in\unsure{\tau|T}}\omega_{S_i}=\unsure{\mathbb{1}|\underline{1}}$ équivaut à la donnée d'un ens à deux éléments $\omega$ (la ``valeur commune'' des $\omega_{S_i}$ identifiés par le syst. transitif d'iso $\bigwedge_{i\in\tau}\omega_{S_i}\simeq\mathbb{1}$) et une famille de \uline{torseurs} $S_i$ $(i\in\tau)$ sous $\mathbb Z/3\mathbb Z\wedge_{\pm1}\omega=\Gamma_\omega$

% L6–10
Considérons alors le \inserted{$\Gamma_\omega$-}torseur :
\[
\deleted{\unsure{\lambda}}T=\deleted{\prod_{i\in\tau}}\bigwedge_{i\in\tau}S_i
\]
alors $G$ opère sur $(\Gamma_\omega,T)$ je dis que $G'$ y opère trivialement. En effet, un antipodisme $\sigma$ induisant $S_i\overset{u}{\simeq}S_j$ et $S_j\overset{u^{-1}}{\simeq}S_i$ induit
\[
S_i\wedge S_j\xrightarrow{\sim}S_j\wedge S_i
\overset{\mathrm{sym}}{\simeq} S_i\wedge S_j
\qquad (x\wedge y)\longrightarrow(u^{-1}y\wedge ux)
\]

% L11–15
\deleted{Or si \unsure{$u'x$|$ux$}} Fixons $a\in S_i$, posons $ua=b\in S_j$, on a donc pour $\gamma,\gamma'\in\Gamma_\omega$\qquad $x=\gamma a,y=\gamma'b$
\[
\sigma(\underbrace{\gamma a\wedge\gamma'b}_{\gamma\gamma'\,a\wedge b})
=(\underbrace{\deleted{\unsure{v|u}}u^{-1}\gamma'b}_{\gamma'a})\wedge
\underset{\gamma b}{(u\gamma a)}
=\deleted{\unsure{\gamma\gamma'}}\underset{\gamma'\gamma\,a\wedge b}{\gamma'a\wedge\gamma b}
\]

% L16–21
Ainsi $G/G'$ opère sur \deleted{\unsure{$T,\omega$}}$(\Gamma_\omega,T)$. Mais $T$ se réalise \inserted{\unsure{canoniqt|canoniquement|canonique}} comme l'ens des \inserted{$3$} autom involutifs $\unsure{\alpha|\deleted{\alpha_\psi}}$ de $S=\coprod_{i\in\tau}S_i$ tels que $\alpha$ induise la transposition de $\tau$ \uline{et} de $\omega$, et on vérifie que $\alpha_t$ conjugue $\alpha_{\unsure{t'|j'}}$ en $\alpha_{\unsure{t''|j''}}$ (en fait, on l'a déjà fait \ldots)
% archive sheet number: 6

